\begin{tabularx}{\textwidth}{|p{3.8cm}|Z|Z|}
\caption{Сравнительный анализ выбранных технологий и альтернатив}
\label{tab:tech_stack_compare} \\

\hline
\multicolumn{1}{|c|}{Технология} &
\multicolumn{1}{c|}{Причины выбора} &
\multicolumn{1}{c|}{Альтернативы и ограничения} \\ \hline
\multicolumn{1}{|c|}{1} & 
\multicolumn{1}{c|}{2} &
\multicolumn{1}{c|}{3} \\ \hline
\endfirsthead

%\multicolumn{3}{r}{} \\[-0.3em]
\multicolumn{3}{r}{{\fontsize{14pt}{14pt}\selectfont Продолжение таблицы \thetable}} \\[0.5em]
\hline
%\multicolumn{1}{|c|}{Технология} &
%\multicolumn{1}{c|}{Причины выбора} &
%\multicolumn{1}{c|}{Альтернативы и ограничения} \\ \hline
\multicolumn{1}{|c|}{1} & 
\multicolumn{1}{c|}{2} &
\multicolumn{1}{c|}{3} \\ \hline
\endhead

uvicorn &
Лёгкий и производительный ASGI-сервер; простота запуска и интеграции с FastAPI; минимальная операционная сложность &
gunicorn -- более сложная конфигурация и ориентирован на production-оркестрацию; hypercorn -- менее распространён и имеет меньшую экосистему \\

\hline

cryptography &
Production-уровень реализации криптографических примитивов; поддержка AEAD-алгоритмов (AES-GCM, ChaCha20-Poly1305); использование OpenSSL; высокая надёжность и широкое применение в индустрии &
PyNaCl -- проще, но ограниченный набор алгоритмов; низкоуровневый OpenSSL API -- более гибкий, но существенно увеличивает риск ошибок реализации \\

\hline

FastAPI &
Нативная поддержка асинхронности (ASGI); высокая производительность; автоматическая генерация OpenAPI; тесная интеграция с Pydantic и строгая типизация &
Flask -- менее строгая типизация, требуется ручная валидация; Django REST Framework -- избыточен для лёгкого сервиса, выше порог сложности \\

\hline

Pydantic v2 + pydantic-settings &
Строгая валидация данных на основе аннотаций типов; высокая производительность (v2); единая модель данных для CLI, API и конфигурации; снижение количества ошибок &
marshmallow -- более многословен и менее удобен; dataclasses -- не предоставляют встроенной валидации и требуют ручной реализации \\

\hline

SQLAlchemy 2.x &
Зрелый ORM с гибкой архитектурой; поддержка как ORM, так и Core; контроль транзакций; независимость от фреймворка &
raw SQL -- больше контроля, но выше сложность и риск ошибок; Django ORM -- привязан к Django и избыточен для данного проекта \\

\hline

pytest &
Лаконичный синтаксис тестов; мощная система фикстур; удобство модульного и интеграционного тестирования; де-факто стандарт в Python &
unittest -- более громоздкий синтаксис и меньшая гибкость при организации тестов \\

\hline

\end{tabularx}	